% Réaction entre deux couples acido-basiques
\begin{donnees}
    Couple acido-basique~:
    \begin{itemize}
      \item $\ce{CH3COOH_{(aq)}}\ttslash{}\ce{CH3COO^{-}_{(aq)}}$
      \item $\ce{NH^{+}_{4(aq)}}\ttslash{}\ce{NH3_{(aq)}}$
      \item $\ce{H3O^{+}_{(aq)}}\ttslash{}\ce{H2O_{(\ell)}}$
      \item $\ce{H2O_{(\ell)}}\ttslash{}\ce{HO^{-}_{(aq)}}$
    \end{itemize}
\end{donnees}

\begin{enumerate}
    \item D'après les données ci-dessus, comment peut-on qualifier l'eau~?
          pourquoi~? 
    \item L'acide acétique réagit-il avec l'eau~? Écrire l'équation de la
          réaction.
    
          Que peut-on dire du pH de la solution obtenue~?
    \item L'ammoniac réagit-il avec l'eau~? Écrire l'équation de la réaction.

          Que peut-on dire du pH de la solution obtenue~?
    \item On dispose de solutions aqueuses d'acide acétique, d'acétate de
          sodium, d'ammoniac et de chlorure d'ammonium.

          Comment représente-t-on ces solutions~? 
    \item Laquelle de ces solutions va réagir avec l'ammoniac~? pourquoi~?

          Écrire l'équation de la solution.
    \item Laquelle de ces solutions réagit avec la solution d'acétate de
          sodium~? pourquoi~?

          Écrire l'équation de la réaction.
\end{enumerate}

